\scanpage{216}
\begin{center}
\uuline{POINTS DE GORENSTEIN D'UN MORPHISME}

par Manuela STOIA
\end{center}

Rappelons qu'un anneau noethérien est de Gorenstein s'il est de dimension injective finie (3). On va utiliser le critère suivant: ``\uline{Un anneau local noetherien $\underline A$, quotient d'un anneau local noethérien régulier $\underline B$, est de Gorenstein si et seulement si $\operatorname{Ext}^{q}_{\underline B}(\underline A,\underline B)=0$ pour tous les entiers $\underline q$ à l'exception d'un $q_{\underline o}$ et $\operatorname{Ext}^{q_o}_{\underline B}(\underline A,\underline B)\simeq\underline A$}'' [(3), prop. 1.3, ch 1, § 1.4].

Soit $\underline k$ un corps algébriquement clos. On va considérer des variétés algébriques $\underline X$ sur $\underline k$ au sens de (4), mais pouvant avoir des éléments nilpotents. On dit que $\underline X$ est de Gorenstein au point $\underline x$ (réspect. de Gorenstein) si l'anneau local $\textcircled{$\mathcal O$}_{\underline X,\underline x}$ est de Gorenstein (réspect. si tous les anneaux $\textcircled{$\mathcal O$}_{\underline X,\underline x}$ sont de Gorenstein).

\uline{Théorème}: \uline{Si $\underline f:\underline X\longrightarrow\underline Y$ est un morphisme plat de variétés algébriques, alors l'ensemble} :
\[
\underline{\underline G_{\underline f}(\underline{\unsure{X|x}})=\{\underline x\in\underline X/\underline X_{\underline f(\underline x)}\quad\text{soit de Gorenstein au point }\underline x\}}
\]
\uline{est ouvert dans $\underline X$ (pour $\underline y\in\underline Y$, on note comme d'habitude par $\underline X_{\underline y}$ la variété algébrique- fibre de $\underline f$ au-dessus de $\underline y$)}.

\uline{Démonstration}. L'énoncé est local sur $\underline X$, donc on peut supposer $\underline X$ affine et soit $\underline j:\underline X\longrightarrow\underline{k^n}$ une immersion fermée. Notons par $\underline i$ l'immersion donné par la composition
\diagram{\begin{tikzcd}[ampersand replacement=\&]
\underline X \arrow[r,"\underline{\textcircled{$\Gamma$}_f}"] \& \underline X\mathbin{x}\underline Y \arrow[r,"\underline j\mathbin{x}\underline{1_Y}"] \& \underline{k^n}\mathbin{x}\underline Y
\end{tikzcd}}
et par $\textcircled{$\pi$}:\underline{k^n}\mathbin{x}\underline Y\longrightarrow\underline Y$ la projection. On a la relation $\underline f=\textcircled{$\pi$}\circ\underline i$. Pour la commodité,
\margin{215}
